Signal Temporal Logic formalizes temporal properties over real-valued signals and provides a quantitative robustness semantics \cite{maler2004,donze2010}. xSTL and AMT combine logic and timed regular expressions for qualitative and quantitative trace analysis \cite{nickovic2018}. \emph{Quantitative Regular Expressions} (QREs) describe modular numerical transformations over streams and have been used to express detectors in both the time and wavelet domains; StreamQRE compiles such specifications modularly for streaming data \cite{abbas2019,mamouras2017}. StreamQL provides an algebra of stream transformations combining relational constructs, dataflow, and temporal operators \cite{kong2020}. MorphoQL differs from these general languages through an event ontology materialized from time--scale maxima and maximum lines, exposed relationally with coefficient-level provenance.

Within data systems, SASE and \code{MATCH\_RECOGNIZE} recognize ordered event sequences \cite{gyllstrom2007,petkovic2022}. T-ReX extends \code{MATCH\_RECOGNIZE} with segment variables and an optimizer designed for time-series patterns \cite{huang2023}. MorphoQL can be viewed as a domain layer upstream of such engines: the extractor creates events, and the compiler can then produce joins, an automaton, or a row-pattern expression.

\subsection{Symbolic Representations and Subsequence Search}

Subsequence search based on Fourier representations, SAX, and the Matrix Profile provides effective solutions for query-by-example, motif discovery, and discord detection \cite{faloutsos1994,lin2007,yeh2016}. SAXRegEx combines a multivariate symbolic representation, regular expressions, and query expansion supporting duration variation and inter-channel offsets \cite{yu2023}. Shape expressions use regular expressions over parametric shapes with a noise-aware semantics \cite{nickovic2021}. MorphoQL mainly differs in its unit of indexing: multiscale events rather than complete numerical windows or segments.

\subsection{Wavelet Maxima and Maximum Lines}

Modulus maxima of a wavelet transform can localize singularities and track their propagation across scales \cite{mallat1992}. Ridges of analytic transforms instead characterize oscillatory components and their instantaneous frequencies \cite{lilly2010,lilly2017}. The present system uses \emph{modulus-maximum lines} for transitions. The term \emph{ridge} is reserved for non-evaluated oscillatory extensions.

\subsection{Natural Language and Time Series}

CLaSP, TS-CLIP, and LaSTR align text with series or segments in learned latent spaces \cite{ito2024,chen2025,dohi2026}. Sonar-TS retrieves candidates before verifying programs on the signal \cite{tan2026}. Grammar of the Wave composes primitives into logical trees and grounds them with vision--language agents \cite{wan2026}; ShapeTalk translates text and sketches into editable constraints \cite{sun2026}. T2SP deterministically transforms series into structured programs of trends, periods, and salient events for language-model reasoning \cite{kim2026}. MorphoQL follows a complementary architecture: natural language may compile into a visible program, but satisfaction of a structured query does not depend on a language model and is executed over an event relation with provenance.

\subsection{Precise Positioning}

\begin{table}[H]
\centering
\caption{Summary positioning. The claimed contribution concerns the interface between a WTMM event relation, compilation, and provenance.}
\small
\begin{tabularx}{\textwidth}{@{}p{3.3cm}p{3.0cm}p{3.0cm}X@{}}
\toprule
Family & Representation & Composition & Difference from MorphoQL \\
\midrule
SDL / SSTS / TSseek & alphabet or segments & regular expressions & no WTMM persistence as a relational attribute \\
QRE / StreamQRE & quantitative stream transformations & modular composition & general language; no materialized WTMM relation with coefficient provenance \\
SAXRegEx & multivariate SAX symbols & regular expressions and expansion & symbolic search over segments rather than a time--scale event relation \\
ShapeSearch & perceptual trends and shapes & algebra, sketch, text & no multiscale event relation with coefficient provenance \\
STL / xSTL & real-valued signal & timed logic & semantics over the signal rather than an extracted relation \\
\code{MATCH\_RECOGNIZE} / T-ReX & rows or segments & regular patterns & morphological events must be defined upstream \\
Shape expressions & parametric shapes & regular expressions & primitive fitting rather than time--scale lines \\
CLaSP / LaSTR / Sonar-TS / T2SP & latent space, features, or program & text / programs & free-form interface or reasoning without an audited relational WTMM representation \\
\textbf{MorphoQL Core 0.2} & signed WTMM events & chain and bounded gaps & multiscale provenance and direct relational compilation \\
\bottomrule
\end{tabularx}
\end{table}

The claim is therefore not ``a new shape language in general,'' but the following:
\begin{quote}
MorphoQL investigates an event relation derived from signed maxima and maximum lines as the substrate of a declarative kernel that can be compiled and audited through time--scale provenance.
\end{quote}

\section{Data Model and Two Levels of Correctness}

\subsection{Series, Extractor, and Event Relation}

Let a corpus of regularly sampled series be
\begin{equation}
X^{(s)}=\bigl((t_i,x_i)\bigr)_{i=0}^{T_s-1},
\qquad s\in\mathcal{S},
\qquad t_{i+1}-t_i=\Delta t.
\end{equation}
A parameterized extractor $\Phi_{\theta}$ produces an ordered, partitioned relation:
\begin{equation}
\Events=\bigcup_{s\in\mathcal{S}}\Phi_{\theta}(X^{(s)}).
\end{equation}
A transition event is the tuple
\begin{equation}
 e=(s,\mathrm{id},\tau,\sigma,\alpha,\rho,a_{\min},a_{\max},\chi,\Gamma),
\end{equation}
where $s$ is the series identifier, $\mathrm{id}$ a content-addressed identifier, $\tau$ the representative time, $\sigma\in\{-1,+1\}$ the sign, $\alpha$ the strength, $\rho$ the persistence, $[a_{\min},a_{\max}]$ the scale span, $\chi$ the extractor name, and $\Gamma$ the ordered sequence of provenance nodes. The identifier is computed as
\begin{equation}
 \mathrm{id}=\code{evt\_}\,\Vert\,\operatorname{SHA256}\!\left(J(e\setminus\{\mathrm{id}\})\right),
\end{equation}
where $J$ is the strict deterministic-JSON serialization of all attributes, including provenance. This invariant is enforced by the public API: an explicit identifier supplied by the caller is accepted only when it equals this address. The identifier is therefore invariant to the insertion of an unrelated event; two exactly identical contents intentionally receive the same address and are rejected by the key constraint rather than hidden under arbitrary names. A node contains at least
\begin{equation}
 n=(a,t,c,b,k),
\end{equation}
where $a$ is scale, $t$ time, $c$ the normalized signed coefficient, $b$ the extraction branch, and $k$ an optional scale index.

The materialized relation is closed by an explicit contract. For every $e\in\Events$,
\begin{equation}
\begin{aligned}
 &(s,\mathrm{id})\text{ is a unique key and both components are nonempty},\\
 &\tau,\alpha,\rho,a_{\min},a_{\max}\text{ are finite whenever present},\\
 &\alpha\geq0,\qquad 0\leq\rho\leq1,\qquad 0\leq a_{\min}\leq a_{\max},\\
 &\Gamma\neq\varnothing,\qquad \operatorname{sign}(c)=\sigma\text{ for every provenance node}.
\end{aligned}
\end{equation}
An \code{EventRelation} instance is also homogeneous in $\chi$. In materialized Core 0.2, the scale span is mandatory and equals the extrema of $\Gamma$. For the four built-in extractors, preprocessing and extraction configurations must exactly equal the registered objects in \path{src/configuration.py}; this requirement also applies to an empty relation that claims the name of a built-in extractor. An empty external relation may omit a configuration, but it is then assigned the explicit status \code{empty\_unconfigured} and cannot be presented as registered. A materialized third-party extension is accepted as a \emph{declared and versioned} configuration when its \code{schema\_version}, \code{config\_kind}, extractor name, and JSON compatibility are valid. Configuration SHA-256 fingerprints are included in the relation.

The API rejects duplicate keys, explicit identifiers inconsistent with content addressing, duplicate content, \code{NaN} or infinite values, negative strengths, persistence outside $[0,1]$, missing provenance, missing scale spans, and mixtures of extractors. Programmatic query parameters are closed as well: both bounds of every \code{Gap} must be finite and satisfy $0\leq l\leq u$; \code{MIN\_STRENGTH} is finite; the deduplication tolerance is either \code{None} or a finite nonnegative real; a projection cannot be empty; and aliases follow the same regular expression as the text parser. Serialization uses \code{allow\_nan=False}, and the application schema rejects silent conversion from a floating-point value or string to an integer. Built-in configurations are compared through deterministic JSON and their fingerprints, not through permissive Python-dictionary equality: payloads \code{2} and \code{2.0} remain distinct. All three strict engines apply the same validation even when they receive a raw event collection directly. The domain accepted by the implementation therefore matches the domain of the formal semantics and the primary key of the SQLite plan.

Every match lies within one series partition $s$; cross-series joins are forbidden by the denotation, the Python engines, and the compiled SQL.

\subsection{Two Complementary Provenance Layers}

We distinguish:
\begin{description}[style=nextline,leftmargin=1.2cm]
  \item[Extraction provenance] the sequence $\Gamma$ of numerical nodes that produced each event;
  \item[Query provenance] the normalized program, its AST, required and observed gaps, threshold, score, postprocessing order, and identifiers of the selected events.
\end{description}
These layers are materialized in an \code{ExecutionWitness}. Its strict deterministic-JSON serialization and round trip are tested. The object includes the schema version, the implementation version and SHA-256 fingerprint, the engine and strict plan, the fingerprint of the sorted relation, the registered or declared preprocessing and extractor configurations and their fingerprints, and the visible-output policy: total ordering, projection, limit, deduplication tolerance, and quasi-duplicate selection rule. It contains exactly as many events and identifiers as the query contains atoms, and exactly one observed and required gap per \code{THEN} operator. A match witness requires at least one visible output and a rank within that output.

Four API levels are separated:
\begin{itemize}
  \item \code{ExecutionWitness.from\_json} performs strict schema deserialization without silent coercion;
  \item \code{verify\_witness\_internal} recomputes the record's internal consistency;
  \item \code{verify\_witness\_against\_relation} recompiles the query, verifies the input fingerprint, and recomputes the strict set through the declared engine;
  \item \code{verify\_witness} reproduces the complete visible pipeline: normative ordering, deduplication, limit, and projection, then compares manifests, the projected row, and the rank.
\end{itemize}
Caller-supplied lists are objects to be compared; they do not become the source of truth. An externally expected policy may be bound explicitly to distinguish intrinsic consistency from authenticity of a claimed protocol. Built-in configurations come from the same \path{src/configuration.py} registry consumed by the extractors and protocol.

\subsection{Execution Witness, Not an Absolute Proof}

Two claims remain distinct:
\begin{enumerate}
  \item $\Events\models q$ means that an event tuple satisfies the syntax and constraints of $q$;
  \item $X^{(s)}$ actually contains a given transition or shape is an empirical claim about the fidelity of $\Phi_{\theta}$.
\end{enumerate}
